%------------------------------------------------------------------------------------
%  SECTION 4.1 : LES RÈGLES D'OR DE LA TYPOGRAPHIE TABULAIRE AVEC BOOKTABS
%------------------------------------------------------------------------------------
\section{Les Règles d'Or de la Typographie Tabulaire et \textsf{booktabs}}
\label{sec:regles_booktabs}

Un tableau n'est pas une grille de tableur : c'est un dispositif typographique destiné à ordonner des données pour faciliter la comparaison intellectuelle.
Dans son manuel de référence devenu le standard mondial des revues à comité de lecture~\cite{fear2020booktabs}, Simon Fear énonce deux axiomes incontournables :

\begin{boiteSynthese}{Les deux axiomes de Simon Fear (paquet \textsf{booktabs})}
  \begin{enumerate}
    \item \textbf{Ne jamais utiliser de filets verticaux.} Les barres verticales entravent le balayage visuel horizontal de gauche à droite et créent un quadrillage lourd et inutile.
    \item \textbf{Ne jamais doubler les filets horizontaux.} Une structure claire s'obtient par la modulation de l'épaisseur des traits et l'équilibre des blancs, non par la multiplication de filets parasites.
  \end{enumerate}
\end{boiteSynthese}

\subsection{La hiérarchie des trois règles horizontales}
\label{subsec:trois_regles_horizontales}

Un tableau scientifique standard s'articule autour de trois filets de graisses distinctes :
\begin{itemize}
  \item \texttt{\textbackslash toprule} : Filet supérieur d'ouverture, d'une épaisseur marquée (généralement \qty{0.08}{\centi\meter}).
  \item \texttt{\textbackslash midrule} : Filet intermédiaire plus fin (\qty{0.05}{\centi\meter}), séparant l'en-tête explicatif des données numériques ou textuelles.
  \item \texttt{\textbackslash bottomrule} : Filet de fermeture d'épaisseur identique à \texttt{\textbackslash toprule}.
  \item \texttt{\textbackslash cmidrule(lr)\{i-j\}} : Filet sous-colonne partiel reliant les colonnes $i$ à $j$. L'option \texttt{(lr)} raccourcit subtilement les extrémités gauche et droite pour séparer nettement deux sous-groupes adjacents.
  \item \texttt{\textbackslash addlinespace} : Insère un espace vertical délicat entre deux blocs logiques de données, évitant de surcharger le tableau d'un trait continu.
\end{itemize}

\subsection{Positionnement normé des légendes}
\label{subsec:position_legendes_tableaux}

Contrairement aux figures dont la légende est systématiquement placée \textbf{en dessous} du graphique, la norme internationale stipule que \textbf{la légende d'un tableau doit impérativement être placée au-dessus} du tableau.
En effet, le lecteur doit prendre connaissance du contexte et de la signification des variables avant d'analyser les valeurs de la matrice.

\texttt{ClasseGuillaume} garantit cette convention grâce à l'option KOMA-Script \texttt{captions=tableheading} et au paramétrage du paquet \textsf{caption} :
\begin{verbatim}
\captionsetup[table]{position = above, skip = 8pt}
\end{verbatim}

\subsection{Démonstration comparative : Tableau de synthèse de mesures}
\label{subsec:demo_synthese_mesures}

Le tableau modulaire ci-dessous illustre l'usage conjugué de \textsf{booktabs}, \textsf{tabularx} et \textsf{siunitx} :

%------------------------------------------------------------------------------------
%  TABLEAU 1 : SYNTHÈSE DE MESURES EXPÉRIMENTALES (booktabs + tabularx + siunitx)
%------------------------------------------------------------------------------------
\begin{table}[htbp]
  \centering
  \small
  \caption{Synthèse des mesures thermomécaniques et validation expérimentale}
  \label{tab:synthese_mesures}
  \begin{tabularx}{\textwidth}{l c c c S[table-format=3.2] X}
    \toprule
    \textbf{Échantillon} & \textbf{Température} & \textbf{Pression} & \textbf{Vitesse} & {\textbf{Rendement (\%)}} & \textbf{Observations qualitatives} \\
    \midrule
    ECH-01 & \qty{22.4}{\degreeCelsius} & \qty{101.3}{\kilo\pascal} & \qty{1.45}{\meter\per\second} & 88.50 & Écoulement laminaire stable sans turbulence. \\
    ECH-02 & \qty{45.0}{\degreeCelsius} & \qty{150.2}{\kilo\pascal} & \qty{2.80}{\meter\per\second} & 92.15 & Amorçage de vortex en sortie de tuyère. \\
    ECH-03 & \qty{75.8}{\degreeCelsius} & \qty{210.0}{\kilo\pascal} & \qty{4.12}{\meter\per\second} & 95.40 & Régime turbulent pleinement développé. \\
    \addlinespace
    ECH-04 & \qty{90.2}{\degreeCelsius} & \qty{280.5}{\kilo\pascal} & \qty{5.60}{\meter\per\second} & 94.80 & Perte de charge par friction pariétale. \\
    ECH-05 & \qty{110.0}{\degreeCelsius} & \qty{350.0}{\kilo\pascal} & \qty{7.25}{\meter\per\second} & 91.20 & cavitation détectée au col d'étranglement. \\
    \bottomrule
  \end{tabularx}
\end{table}



Comme le met en évidence le~\cref{tab:synthese_mesures}, l'absence de traits verticaux allège considérablement la lecture et confère une élégance immédiate à la présentation des résultats.

